\documentclass[a4paper,12pt]{article}
\usepackage{graphicx}
\usepackage[a4paper, top=1cm, bottom=2cm, left=2cm, right=2cm]{geometry}
\usepackage{xcolor}
\usepackage{hyperref}
\usepackage{eso-pic}

\usepackage[french]{babel}
\usepackage[T1]{fontenc}
\usepackage{mathrsfs}
\usepackage{amsmath}
\usepackage{amsfonts}
\usepackage{amssymb}

\usepackage{tikz}
\usetikzlibrary{arrows, shapes.geometric, fit}

\newcounter{correction}

\newcommand{\question}{%
  \refstepcounter{correction}%
  \textbf{\textcolor{black}{Question \thecorrection :}} \ignorespaces
}

\AddToShipoutPicture{
    \AtTextCenter{%
        \makebox[0pt]{\rotatebox{80}{\textcolor[gray]{0.6}{\fontsize{10cm}{10cm}\selectfont PGB}}}
    }
}

\begin{document}
\hrule
\begin{center}
    \begin{tabular}{@{} p{5cm} p{5cm} p{5cm} @{}}
        Lycée Dindéfélo & \quad\quad\quad Test 3 & 21 Octobre 2026 \\
    \end{tabular}
    \\[-0.01cm]
    \hrule
\end{center}
\begin{center}
    \textbf{\Large Sommes Algébriques et Puissances de Radicaux} \\[0.2cm]
    \textbf{\large Professeur : M. BA} \\[0.2cm]
    \textbf{Classe : 3 ème} \\[0.2cm]
    \textbf{\small Durée : 15 minutes} \\[0.2cm]
    \textbf{\small Note :\quad\quad /5}
\end{center}

\textbf{\small Nom de l'élève :} \underline{\hspace{8cm}} \\[0.5cm]

\question (1,5 point) Écrire les radicaux suivants sous la forme $a\sqrt{b}$ avec $a \in \mathbb{N}$ et $b \in \mathbb{N}$ le plus petit possible :
\[
\begin{aligned}
\sqrt{27}&=\sqrt{\textcolor{red}{3^2} \times 3}\\
&=\textcolor{red}{\sqrt{3^2}} \times \textcolor{red}{\sqrt{3}}\\
&=\textcolor{red}{3\sqrt{3}}\\
\end{aligned}
\quad\quad\quad\quad\quad
\begin{aligned}
\sqrt{72}&=\sqrt{\textcolor{red}{2^2\times3^2} \times 2}\\
&=\textcolor{red}{\sqrt{2^2}} \times \textcolor{red}{\sqrt{3^2}} \times \textcolor{red}{\sqrt{2}}\\
&=\textcolor{red}{2} \times \textcolor{red}{3} \times \textcolor{red}{\sqrt{2}}\\
&=\textcolor{red}{6\sqrt{2}}\\
\end{aligned}
\quad\quad\quad\quad\quad
\begin{aligned}
\sqrt{75}&=\sqrt{\textcolor{red}{5^2} \times 3}\\
&=\textcolor{red}{\sqrt{5^2}} \times \textcolor{red}{\sqrt{3}}\\
&=\textcolor{red}{5\sqrt{3}}\\
\end{aligned}
\]

\vspace{0.8cm}

\question (2 points) Simplifier les sommes algébriques suivantes :
\[
\begin{aligned}
A&= 2\sqrt{27} - \sqrt{75}\\
&=\textcolor{red}{2(3\sqrt{3}) - 5\sqrt{3}}\\
&=\textcolor{red}{6\sqrt{3} - 5\sqrt{3}}\\
&=\textcolor{red}{(6-5)\sqrt{3}}\\
& =\textcolor{red}{ \sqrt{3} }	
\end{aligned}
\quad\quad\quad\quad\quad\quad\quad\quad
\begin{aligned}
B&= \sqrt{72} + 3\sqrt{18} - \sqrt{50}\\
&=\textcolor{red}{6\sqrt{2} + 3(3\sqrt{2}) - 5\sqrt{2}}\\
&=\textcolor{red}{6\sqrt{2} + 9\sqrt{2} - 5\sqrt{2}}\\
&=\textcolor{red}{(6 + 9 - 5)\sqrt{2}}\\
&=\textcolor{red}{ 10\sqrt{2} }
\end{aligned}
\]

\vspace{0.8cm}

\question (1,5 point) Simplifier les expressions suivantes comportant des puissances sous le radical :
\[
\begin{aligned}
C&=\sqrt{5^6}\\
&=\textcolor{red}{\sqrt{(5^3)^2}}\\
&=\textcolor{red}{5^3}\\
&=\textcolor{red}{125}\\
\end{aligned}
\quad\quad\quad\quad\quad\quad\quad\quad
\begin{aligned}
D&=\sqrt{2^7}\\
&=\textcolor{red}{\sqrt{2^{3\times{2+1}}}}\\
&=\textcolor{red}{2^3\sqrt{2}}\\
&=\textcolor{red}{8\sqrt{2}}\\
\end{aligned}
\]

\end{document}